\section{Arquitetura Python}

O arquivo \codefile{main.py} é o ponto de entrada do simulador. Ele carrega a configuração, cria a fonte de impulso, instancia o modelo de bobina, resolve as equações diferenciais, salva CSVs e gera figuras e GIFs. O mesmo script executa quatro cenários: modelo Pi, modelo T, baixa capacitância e alta capacitância. Essa escolha torna o projeto útil não apenas para produzir um caso único, mas para comparar como a topologia e a capacitância afetam a propagação do surto.

\codefile{src/models/distributed\_coil.py} concentra a formulação elétrica. Para o modelo Pi, o estado é composto pelas tensões dos nós internos e de saída, além das correntes de cada ramo série. As derivadas são calculadas por KCL nos capacitores e KVL nos ramos R--L. Para o modelo T, cada seção possui um nó intermediário capacitivo e correntes de junção; isso desloca os pontos de observação para o meio das seções e explica por que os vetores de posição do modelo T possuem mais pontos.

\codefile{src/sources/impulse\_source.py} define a fonte de tensão. A implementação oferece a onda dupla exponencial, usada no caso base, e uma alternativa rampa-exponencial. \codefile{src/solvers/time\_domain\_solver.py} usa \codefile{scipy.integrate.solve\_ivp} com o método RK45. A documentação do SciPy descreve esse método como uma rotina para problemas de valor inicial, com controle adaptativo de erro por fórmulas Runge--Kutta embutidas \cite{scipy_solve_ivp}.

\codefile{src/utils/result\_processor.py} calcula grandezas derivadas importantes para engenharia: pico de tensão por nó, maior diferença de tensão entre nós adjacentes, tensão de entrada, tensão de saída e relação de transferência. \codefile{src/visualization/plot\_generator.py} transforma esses dados em PNGs estáticos. \codefile{src/visualization/gif\_generator.py} gera animações da onda de tensão, mapas de calor animados e comparações entre cenários.

Três scripts auxiliares fazem a ponte com o caso ATP e com visualizações adicionais. \codefile{run\_atp.py} executa o ATP, lê o arquivo binário \codefile{.pl4} e gera gráficos Plotly com os \emph{resultados reais do ATP}. \codefile{scripts/plot\_plotly.py} gera as mesmas visualizações Plotly a partir do \emph{simulador Python}, reutilizando o pacote \codefile{src/} e a configuração JSON — ele substitui os antigos \codefile{simulate\_direct.py} e \codefile{save\_images.py}, que duplicavam o modelo físico com parâmetros próprios (risco apontado na auditoria do projeto). \codefile{scripts/compare\_python\_atp.py} executa a validação cruzada quantitativa entre as duas soluções (Seção~\ref{sec:validacao-quantitativa}). Nenhum dos scripts usa caminhos absolutos: os caminhos derivam da raiz do repositório e o executável do ATP pode ser indicado por argumento de linha de comando ou pela variável de ambiente \codefile{ATP\_EXE}.

\begin{table}[H]
\centering
\caption{Fluxo de execução do simulador Python.}
\label{tab:fluxo-python}
\begin{tabular}{p{0.23\linewidth}p{0.67\linewidth}}
\toprule
Etapa & Função técnica \\
\midrule
Configuração & Lê parâmetros físicos, temporais e de saída em JSON. \\
Fonte & Gera a excitação transitória aplicada ao nó de entrada. \\
Modelo & Monta a rede Pi ou T por seções distribuídas. \\
Solver & Resolve o sistema de EDOs no domínio do tempo. \\
Processamento & Calcula picos, gradientes e relação de transferência. \\
Visualização & Produz PNGs e GIFs para interpretar a onda viajante. \\
\bottomrule
\end{tabular}
\end{table}
